% HD-I23. GPT-6 Astra Pro /root/imprimitive; inherited effort, not numerically exposed.
% 21 November 2026. Author draft. Original coefficient projectors and polarized maps.
\section{HD-I23: the actual rank-four support and field reduction}
\label{sec:rank-four-support-field}

Let \(K\) be one of
\[
 \mathbf Q(i),\quad \mathbf Q(\sqrt{-d}),\qquad
                         d=2,7,11,19,43,67,163.             \tag{I23.1}
\]
Retain the natural-source conventions of HD-S45 and HD-R15: the
finite coefficient group \(G\subset GL_K(V)\) is faithful and
absolutely irreducible, its exact character field is \(K\), and
the prescribed rational primitive Jacobian corner has signature
\((3,3)\), restricted polarization, and period rank nine.
The reflection and bireflection endpoints are those already
supplied in HD-S36, HD-G27, HD-S41--HD-S43 and HD-G30. In this section the actual
group has no nonidentity element of moving rank one, two, or three.
For its sphere-cover source, with \(m=\dim_KV\),
\[
       b\ge12,\qquad \sum_jc_j=2m+6,\qquad c_j\ge4,
                               \qquad m\ge21.              \tag{I23.2}
\]
Set
\[
       N=\langle x\in G:\operatorname{rank}(x-I)=4\rangle.
                                                               \tag{I23.3}
\]
This is an actual normal subgroup of \(G\), since conjugation preserves
the specified moving rank. If it is
trivial, all branch codimensions are at least five and (I23.2)
gives \(m\ge27\). For nontrivial \(N\), the following construction
gives its complete isotypic support reduction. In particular its
constituent multiplicity \(e\) and Galois-orbit size \(h\) obey
\[
 (e,h)\in\{(1,1),(1,2),(1,4),(2,1),(2,2),(4,1)\}.         \tag{I23.4}
\]
The cases \((1,4)\) and \((2,2)\) have period rank at most four.
The case \((4,1)\) has an actual tensor presentation with a
four-dimensional multiplicity factor; outside the supplied sources,
its reflection constituent is the rational standard \(A_r\)
module, \(r\ge6\), or the rational \(E_6,E_7,E_8\) module.
The cases with \(eh=2\) reduce to the specified bireflection
constituents in paragraph7. The pair \((1,1)\) gives an absolutely
irreducible normal subgroup on a possibly smaller actual source;
either that source returns to a preceding moving-rank branch,
or this subgroup is generated by actual rank-four elements.
All replacements below
are accompanied by maps from the originally prescribed corner.

\paragraph{1. Literal isotypic projectors and the free Galois action.}
Over \(\mathbf C\), write the distinct \(N\)-isotypic components
as \(U_i=W_i\otimes\mathbf C^e\), where \(\dim W_i=r\).
Transitivity follows from irreducibility of \(G\); hence \(r,e\)
are independent of \(i\). If \(\chi_i\) is the character of
\(W_i\), the actual matrix projectors are
\[
       p_i=\frac r{|N|}\sum_{n\in N}\chi_i(n^{-1})\rho(n).
                                                               \tag{I23.5}
\]
Character averaging gives \(p_ip_j=\delta_{ij}p_i\),
\(\sum p_i=I\), and image \(U_i\). Equivalently these are
the complete primitive central projectors of \(\mathbf C\rho(N)\);
the averaging matrix units give its literal restriction isomorphism
\(\mathbf C\rho(N)=\bigoplus_i\operatorname{End}(W_i)\otimes I_e\).
They are orthogonal for the original positive invariant coefficient
form, because an off-diagonal Hermitian pairing would give an
intertwiner between inequivalent finite-group representations.

Take a cyclotomic splitting field \(L=K(\zeta_a)\) containing all
characters. Its Galois group permutes exactly the projectors (I23.5),
and this permutation action commutes with that of \(G\), since all
\(\rho(g)\) have entries in \(K\). Let \(\mathcal G\) be the
image of that Galois action. A permutation commuting with a transitive
group and fixing one label fixes every label. Thus \(\mathcal G\)
acts freely on the labels; write \(h=|\mathcal G|\).

For an actual rank-four generator \(x\), put
\(a_i(x)=\operatorname{rank}(x-I|_{W_i})\). The ranks of
\((x-I)p_i\), computed by minors, are unchanged by a Galois
automorphism, so \(a_i(x)\) is constant on each Galois orbit.
Consequently
\[
                  4=e\sum_i a_i(x)\quad\Longrightarrow\quad eh\mid4.
                                                               \tag{I23.6}
\]
This proves (I23.4). Let \(E=L^{\ker(\operatorname{Gal}(L/K)
\to\mathcal G)}\). It is a Galois extension of degree \(h\).
The projectors have entries in \(E\), and its Galois group is
cyclic of order two for \(h=2\), and cyclic of order four or
the four group for \(h=4\). These assertions use the actual
permutation action; no selected block system is assumed to descend.

\paragraph{2. The field action and the exact smaller source.}
Choose one Galois orbit \(\Omega\), set
\(p_\Omega=\sum_{i\in\Omega}p_i\), and put
\(U=p_\Omega V\), an actual \(K\)-space of dimension \(her\).
Let \(H\) be its stabilizer in \(G\). The original representation
is literally \(\operatorname{Ind}_H^G U\): for coset
representatives \(t_j\), the map is
\((u_j)\mapsto\sum_j\rho(t_j)u_j\), with inverse given by
the orthogonal orbit projections. Any proper complex invariant
subspace of \(U\) would induce a proper one of \(V\); thus
the \(H\)-representation is absolutely irreducible.
Its character field is exactly \(K\), since the induced trace
formula gives \(K\subset\mathbf Q(\chi_U(H))\subset K\).

Choose \(p_0\) in \(\Omega\), label
\(p_\tau=\tau(p_0)\), and define the actual field embedding
\[
 \iota:E\longrightarrow\operatorname{End}_K(U),\qquad
                a\longmapsto\sum_{\tau\in\operatorname{Gal}(E/K)}
                                             \tau(a)p_\tau. 
                                                               \tag{I23.7}
\]
The matrices are Galois invariant, multiplication is componentwise,
and a nonzero scalar has no zero component, so this is injective.
The map from this \(E\)-vector space \(U\) to
\(p_0(U\otimes_KE)\) is \(u\mapsto p_0u\); its inverse is
\(w\mapsto\sum_\tau\tau(w)\). The two composites are the
identity by orthogonality of the projectors. Hence \(\dim_EU=er\).
The stabilizer \(H\) acts semilinearly and maps onto
\(\operatorname{Gal}(E/K)\). Indeed its permutation on
\(\Omega\) commutes with the regular Galois action and is
transitive; because that action is abelian, its centralizer is the
same regular group. Conjugating (I23.7) displays the corresponding
field automorphism explicitly.

Let \(\chi\) be the irreducible character on \(W_0\), and
\(F=\mathbf Q(\chi(N))\). On the displayed \(E\)-space
\(\chi(n)=e^{-1}\operatorname{tr}_E(n)\), so \(F\subset E\).
An automorphism of \(E/K\) fixing \(\chi\) would fix its
isotypic projector (I23.5), contrary to freeness. Thus
\[
                                  KF=E.                    \tag{I23.8}
\]
This proves the exact constituent field assertion even when \(h=4\).

Here are the geometric maps for every use of this block replacement.
Let \(\widetilde C\) be the original full-\(G\) curve,
\(H_0=\ker(H\to GL_K(U))\), and set
\[
 Y=\widetilde C/H,
 \quad C'=\widetilde C/H_0,
 \quad \alpha:\widetilde C\longrightarrow C'.             \tag{I23.9}
\]
The induced coefficient functions and inverse are
\(F_v(x)=p_\Omega\rho(x)v\),
\(F(hx)=\sigma(h)F(x)\), and
\(F\mapsto\sum_{H\backslash G}\rho(x)^{-1}F(x)\).
On each branch cycle the fixed-vector equations give the identical
return-monodromy kernel on both sides, so this includes the
middle-extension stalks. For \(v\in U(K)\setminus0\), with
\(q=h(v,v)\), the actual rational corners are equal in
\(\mathbf Q[G]\):
\[
 \frac{2m}{|G|q}\sum_{g\in G}\operatorname{Re}h(v,\rho(g^{-1})v)g
 =\frac{2\dim_KU}{|H|q}\sum_{h\in H}
                         \operatorname{Re}h(v,\sigma(h^{-1})v)h.
                                                               \tag{I23.10}
\]
Coefficients off \(H\) vanish and \(m/|G|=\dim_KU/|H|\).
After clearing denominators, pullback \(\alpha^*\) is a
surjective isogeny from the corresponding connected primitive
factor of \(JC'\) onto this corner in \(J\widetilde C\), and
\[
 \alpha_*\alpha^*=[|H_0|],\qquad
 (\alpha^*)^*\theta_{J\widetilde C}=|H_0|\theta_{JC'}.
                                                               \tag{I23.11}
\]
For an original line \(Kw\) outside \(U\), use
\(T(x)=v h(w,x)/h(w,w)\) and
\[
 b=\frac{2m}{|G|}\sum_g\operatorname{Re}\operatorname{tr}
                         (T\rho(g^{-1}))g,
 \qquad b^\dagger b=\frac{h(v,v)}{h(w,w)}e_w.                \tag{I23.12}
\]
The matrix-coefficient averaging identities give these equations
on the specified block and its conjugate, and zero on all others.
An integer-cleared \(b\) is an actual isogeny with multiplier
that positive rational ratio times an integer square. On a
six-dimensional \(K\)-Hermitian space the multiplier changes
the determinant by a sixth power, a \(K/\mathbf Q\) norm.
Thus the ordinary determinant class, actual \(K\)-action,
and period rank are retained. This proves the original-corner
comparison for every subsequent descended block as well.
For any such coefficient source of rank at least four, period rank
nine forces its base to have genus zero: genus at least two would
contribute at least eight to its parabolic dimension six, and in
genus one the marked-base dimension is at most six. Thus all later
uses of the sphere budget on a rank-at-least-four replacement are
justified by the actual quotient source.

\paragraph{3. Rank-one spectral projectors split the reflection factor.}
Suppose \(eh=4\). Equation (I23.6) says that each generator is
supported on precisely one Galois orbit and is a reflection on
each of its irreducible constituents. Subgroups supported on
distinct orbits commute and intersect trivially. On the selected
orbit their actual image \(B\) is therefore finite and generated
by reflections on its irreducible constituent \(W\).

Let \(\mathcal A=E[B]\subset\operatorname{End}_E(U)\).
Its complexification on the chosen embedding is
\(M_r(\mathbf C)\otimes I_e\); the ordinary irreducible
character is \(e^{-1}\operatorname{tr}_E\), valued in \(E\).
For a reflection \(x\), its nonidentity eigenvalue is
\(\lambda=\chi(x)-r+1\in E\). The literal idempotent
\[
 P=\frac{x-I}{\lambda-1}\in\mathcal A,
 \qquad P\mathcal AP=EP,\quad
 \dim_E\mathcal AP=r,\quad \dim_EPU=e                         \tag{I23.13}
\]
splits this original algebra. To specify the maps, choose
\(a_j,b_j\in\mathcal A\) with \(\sum_j a_jPb_j=1\),
possible because \(\mathcal AP\mathcal A\) becomes the whole
matrix algebra after scalar extension. Then
\[
 \mathcal AP\otimes_EPU\longrightarrow U,\quad aP\otimes u\mapsto au,
 \qquad u\longmapsto\sum_j a_jP\otimes Pb_ju                 \tag{I23.14}
\]
are inverse maps: on a pure tensor use
\(Pb_jaP\in EP\) and the displayed identity. Left multiplication
identifies \(\mathcal A\) with
\(\operatorname{End}_E(\mathcal AP)\). This constructs the
actual constituent and multiplicity spaces over \(E\); no
Schur-index or tensor-splitting assumption is needed.
The original \(H\) normalizes this algebra and is semilinear
over \(E/K\). On (I23.14) it acts by semilinear tensor maps.
For an \(E\)-linear element this follows by sending a matrix-unit
basis to its conjugate: choose a nonzero vector in the first
image line and transport it by the conjugated matrix units to
obtain an intertwining matrix over \(E\); its remaining
commutant is \(I_r\otimes\operatorname{End}_E(PU)\).
Whenever such a tensor has finite order \(n\), each factor
is semisimple: \(A^n\otimes B^n=I\) implies
\(A^n=aI\), \(B^n=a^{-1}I\), and the two scalar-power
polynomials have distinct roots in characteristic zero.

We use Cohen's published classification in the following exact
range: an irreducible complex reflection group of rank at least
five is either imprimitive \(G(a,b,r)\), a rational real Weyl
group, or \(G_{33}\) of rank five or \(G_{34}\) of rank six,
both of exact character field \(\mathbf Q(\omega)\).
The complete classification is Cohen, \emph{Finite complex
reflection groups}, Ann. Sci. ENS (4) \textbf{9} (1976),
379--436, Theorems2.4 and5.12, the real classification4.4(i),
the explicit rank-five/six roots on p.409 and Table4.16 on p.412.
The exact fields from these matrices are recorded in HD-I11
and HD-G30. The rational irreducible primitive real groups in
rank at least six are \(A_r,E_6,E_7,E_8\); \(B_r,D_r\)
are imprimitive.

\paragraph{4. The case \((e,h)=(1,4)\).}
Here \(U\) is the underlying \(K\)-space of \(E^r\).
For \(r=1\), any nonidentity \(E\)-linear scalar fixes zero,
and a nontrivial semilinear automorphism fixes a \(K\)-space
of dimension at most two. To see the latter bound, divide two
nonzero fixed vectors to obtain an element of \(E^\tau\);
the fixed space is at most one-dimensional over \(E^\tau\).
Thus every nonidentity coefficient has moving rank at least two
on the rank-four source, whose period rank is at most four by
paragraph8 below.

For \(r\ge2\), a nonidentity \(E\)-linear operator has
\(K\)-moving rank at least four. A nontrivial field automorphism
permutes the four complex embeddings freely in cycles of length
two or four; its fixed dimension is at most half the total,
so its moving rank is at least \(2r\ge4\). A rank-nine
source therefore has \(4r\ge21\), hence \(r\ge6\).
The field equality (I23.8), \([E:K]=4\), excludes every
primitive reflection group in the preceding list. Thus \(B\)
is imprimitive.

Here is the actual field and projector descent for that conclusion.
In the standard \(G(a,b,r)\) model let \(\zeta\) be its
primitive \(a\)-th root. For \(a>2\), diagonal elements with
entries \((\zeta,\zeta^{-1},1,\ldots)\) and
\((\zeta,\zeta,\zeta^{-2},1,\ldots)\) give character values
\(c=\zeta+\zeta^{-1}\) and \(v=2\zeta+\zeta^{-2}\), after
subtracting rational constants, and
\[
                   \zeta=\frac{v-c^2+1}{2-c}.               \tag{I23.15}
\]
Thus its exact character field is \(\mathbf Q(\zeta)\), and
\(E=K(\zeta)\). The case \(a\le2\) has rational character
and is already excluded. Solve the finite intertwiner equations
over \(E\) between the actual module (I23.14) and these matrices;
their complex solution space is one-dimensional and contains an
isomorphism, so every nonzero \(E\)-solution is invertible.
This is an actual \(E\)-linear comparison with the standard model.

The diagonal subgroup \(D\) is characteristic in \(B\): it
is abelian and normal, and every abelian normal subgroup maps
trivially to \(S_r\), \(r\ge5\). Its coordinate characters
are distinct. Their averaging projectors are rank-one matrices
over \(E\), are permuted by every semilinear normalizer of
\(B\), and transfer through (I23.7) to actual \(K\)-rank-four
blocks of \(U\). They are orthogonal, by the original invariant
form and distinct diagonal characters. Each block stabilizer
acts on the underlying \(K\)-space of \(E\) by
\(u\mapsto a\tau(u)\); the fixed-space argument above shows
that every nonidentity element has moving rank at least two.
The actual quotient and corner maps (I23.9)--(I23.12) reduce
the period rank to at most four.

\paragraph{5. The case \((e,h)=(2,2)\).}
Now \(U=E^r\otimes_E E^2\) as in (I23.14), of \(K\)-rank
\(4r\). For \(r=1\), nonidentity \(E\)-linear operators have
\(K\)-moving rank at least two, and a semilinear field involution
exchanges two complex two-spaces, hence also has moving rank at
least two. The rank-four period bound again applies.

For \(r\ge2\), an \(E\)-linear tensor \(A\otimes B\)
with both factors nonscalar has moving rank at least \(r\)
over \(E\). This is the eigenvalue multiplicity inequality
\(\dim\ker(A\otimes B-I)\le r\) for a nonscalar
two-dimensional factor. If only the first factor is active its
moving rank is a positive multiple of two; if only the second
is active it is at least \(r\). Thus every nonidentity
\(E\)-linear element has \(K\)-moving rank at least four.
A nontrivial field involution exchanges two complex \(2r\)-spaces,
so has moving rank at least \(2r\ge4\).

The map of the actual source group onto \(\operatorname{Gal}(E/K)\)
is surjective. A generating product-one sphere tuple therefore
has at least two entries in its nontrivial coset. With \(b\ge12\),
its budget gives
\[
                    8r+6\ge4r+4(b-2)\ge4r+40,
                    \qquad r\ge9.                         \tag{I23.16}
\]
Primitive reflection groups of rank at least nine are rational
real groups, contrary to \(KF=E\ne K\). Thus \(B\) is
imprimitive. Apply (I23.15) and the actual \(E\)-intertwiner
as in paragraph4. The characteristic diagonal projectors on
its first factor, tensored with \(I_2\), give actual
\(K\)-rank-four blocks, each the underlying space of \(E^2\).
Its stabilizer is \(E/K\)-semilinear; a nonidentity linear
element has even positive \(K\)-moving rank, and a field
involution has moving rank at least two. The quotient source
therefore again has period rank at most four.

\paragraph{6. The case \((e,h)=(4,1)\) and its exact tensor endpoint.}
Equation (I23.14) is now the actual \(K\)-linear map
\[
                  U=W_K\otimes_K M_K,\qquad \dim_K M_K=4,
                                                               \tag{I23.17}
\]
with \(B\) acting only on \(W_K\). Constituents of rank
\(r\le3\) give an actual coefficient source of rank at most
twelve, already covered by HD-S40 and its seven-field extension.
For \(r\ge4\), both nonscalar tensor factors have moving rank
at least \(r+4-2\); this follows by maximizing the matching
eigenvalue multiplicities, which gives fixed dimension at most
\((r-1)3+1\). A sole active factor has moving rank at least
\(4\) or \(r\). Thus the rank-nine source has
\(4r\ge21\), hence \(r\ge6\).

The classification and \(F\subset K\) exclude \(G_{34}\).
In the imprimitive case its exact diagonal-character projectors
are defined over \(K\), by (I23.15), and tensoring them with
\(I_4\) gives actual \(K\)-rank-four blocks preserved by the
full source group. The rank-four source comparison in paragraph8
gives a supplied determinant. The remaining constituent is
\(A_r\), \(r\ge6\), or \(E_6,E_7,E_8\), with its rational
reflection representation. The finite intertwiner equations
between (I23.14) and that rational model have a nonzero
\(K\)-solution, which is invertible by irreducibility.
This specifies the actual \(K\)-basis change, not just the
constituent's trace field.

For completeness the resulting overgroup is also literal.
The normalizer of a simply laced real Weyl group in its
irreducible complex reflection representation is
\(\mathbf C^*W^{\pm}\), where \(W^{\pm}=\langle W,-I\rangle\).
Indeed the unique invariant quadratic form is preserved up to
scalar; normalize that scalar, and each reflection line of norm
two maps to a root or its negative. Here every reflection of the
Weyl group is a root reflection: its fixed hyperplane, if different
from all root hyperplanes, would contain a regular point, whose
chamber would then have a nonidentity stabilizer. The simple
transitivity of chambers excludes that possibility. The resulting matrix is real
on a root basis and is a root-system automorphism. The root
automorphism classification gives \(W\) times its diagram
automorphisms; for \(A_r,E_6\) the extra involution is
\(-w_0\), while \(E_7,E_8\) have none. To check the former
identification, \(-w_0\) preserves the simple-root basis;
\(-I\notin W\) since \(W\) acts trivially on the weight/root
lattice quotient, whereas \(-I\) acts nontrivially on
\(\mathbf Z/(r+1)\) for \(A_r\), and on \(\mathbf Z/3\)
for \(E_6\). These quotients follow from their displayed Cartan
matrices. Thus this is the nontrivial diagram involution in Table1.
The exact published
locator is Humphreys, \emph{Introduction to Lie Algebras and
Representation Theory}, Theorem10.3(a),(b),(e), pp.51--52,
and section12.2, pp.65--66, including Table1.
These original pages were freshly read for the present argument.

Apply the matrix-unit intertwiner from paragraph3 to each actual
\(K\)-linear overgroup element. Its first factor is a scalar
times a rational \(w\in W^{\pm}\); the scalar lies in \(K\)
by comparing any nonzero entry and is absorbed into the second
factor. Thus the actual group is the tensor image of the finite
group
\[
 \widehat H=\{(w,B')\in W^{\pm}\times GL_4(K):w\otimes B'\in H\},
 \qquad\ker(\widehat H\to H)=\{(I,I),(-I,-I)\}.             \tag{I23.18}
\]
It contains \(W\times1\). Its multiplicity projection
\(\Gamma_4\) is finite and absolutely irreducible: a proper
invariant subspace of its four-dimensional module would tensor
with \(W_K\) to give a proper invariant subspace of the
original coefficient representation. Every original trace is
\(\chi_W(w)\chi_{\Gamma_4}(B')\), with \(\chi_W(w)\in\mathbf Q\).
Consequently its exact character field is \(K\), since
\(K\subset\mathbf Q(\chi_{\Gamma_4})\subset K\).
For \(A_r\), absorb the sign into \(B'\); the
augmentation permutation matrices give a unique pair, so
\(H=S_{r+1}\times\Gamma_4\) literally. For the exceptional
types retain (I23.18), including its displayed kernel. These
are the precise tensor-four endpoints for HD-R19; no arbitrary
central extension is substituted for the original cover group.

\paragraph{7. The two remaining multiplicity--field pairs.}
For \((e,h)=(1,2)\) or \((2,1)\), the restrictions of the
rank-four generators to a selected orbit have moving rank two
or four. A rank-two element in the actual smaller source group
is covered by the existing reflection/bireflection-source results
and their literal quotient maps. Otherwise each nonidentity
restriction has moving rank four, and its irreducible
constituent has moving rank two. The projected group on that
constituent is therefore a finite irreducible group generated
by actual bireflections.

Here the constituent rank is at least eleven in every remaining
period-rank-nine case. For \((1,2)\), the source is the underlying
\(K\)-space of \(E^r\). Linear elements have even moving rank;
field-involution elements have moving rank at least \(r\).
For \(r\ge4\) there is no moving rank one or three, and after
the rank-two endpoint every nonidentity element has rank at least
four. For \((2,1)\), the source is complex-isotypic
\(W\otimes\mathbf C^2\); every normalizer is a tensor operator
on these actual complex spaces, by the matrix-unit argument.
Both active factors have moving rank at least \(r\), a sole
second factor at least \(r\), and a sole first factor an even
rank. The same conclusion follows for \(r\ge4\).
The ranks \(r\le3\) give coefficient rank at most six and
are already covered. In the remaining cases the sphere budget
is \(4r+6\ge48\), so \(r\ge11\).
These are the exact bireflection inputs to HD-I25. When
\((e,h)=(1,1)\), the projected normal subgroup is already
absolutely irreducible over \(K\). Its generators have moving
rank at most four. Moving rank one or two returns to the supplied
endpoints; moving rank three returns, through the same actual
maps, to the prior rank-three reduction. Otherwise all nonidentity
generator restrictions have rank four. If there was more than one
isotypic orbit, this is a strict decrease in coefficient rank.

\paragraph{8. The small coefficient bounds used above.}
Every actual block stabilizer used here is absolutely irreducible
and has exact character field \(K\), by the induction argument
in paragraph2. For a rank-four system whose first parabolic
cohomology has dimension six, the exact formula is
\(6=4(2g-2)+\sum c_j\). If each nonidentity \(c_j\ge2\),
genus zero gives at most seven branch points and period rank
at most four; genus one gives at most three points and rank
at most three; larger genus is impossible. Without the
codimension-two restriction, a rank-nine such source must have
genus zero and at least twelve points. Its total codimension
is fourteen, so it contains a reflection. Its determinant is
therefore supplied by HD-S36 for the Gaussian field and
HD-G27 for the seven other fields. In each application the
original corner comparison is (I23.9)--(I23.12), with the
specified polarization multipliers.
